\scanpage{86}
La flèche composée de (11) est alors une flèche $\widetilde s_{i'_n}\longrightarrow\widetilde t_{i'_1}$ qui se trouve dans une catégorie fibre, donc \inserted{à flèches génératrices \unsure{données}} on peut trouver un mot \inserted{$M_{\alpha,i_1}$} en les $a_j$ tels que
\[\deleted{\unsure{M_\alpha}}\]
\diagram{\begin{tikzcd}[ampersand replacement=\&]M_{\alpha,1}((a_j)):\quad\widetilde t_{i'_1}\arrow[r,squiggly]\&\widetilde s_{i'_n}\end{tikzcd}}
et que la composée de \deleted{la flè.} cette flèche avec la séquence précédente soit l'identité :
\[(12)\qquad\boxed{M_{\alpha,1}((a_j))u'_{i'_{\alpha,1}}M_{\alpha,2}((a_j))u'_{i'_{\alpha,2}}\cdots M_{\alpha,\unsure{i'_{\alpha,n_\alpha}}}((a_j))u'_{i'_{\alpha,n_\alpha}}=1}\]
\margin{relations (8) ``remontées''}
Supposons donné de plus une \deleted{\unsure{un}}\inserted{famille} de relations sur les flèches-fibres $a_j$
\[(13)\qquad\boxed{N'_\beta((a_j))=1}\]
telles que celles de ces relations qui sont relatives à une catégorie-fibre fixée $\widetilde{\mathcal T}_s$, soient \unsure{définissantes} pour celle-ci (pour la \deleted{\unsure{aux}}\inserted{famille génératrice} de flèches \deleted{\unsure{génératrices}}\inserted{$(a_j)$}) ! \unsure{Supposons} enfin que pour $\forall\widetilde s\in\widetilde S$, les \deleted{\unsure{flèch.}} flèches $a_j$ qui sont des lacets en $\widetilde s$ forment une
\mspage{85}
